%--------------------------------------------------------------------%
% BAB II ANALISIS DAN PERANCANGAN
%--------------------------------------------------------------------%
\chapter{ANALISIS DAN PERANCANGAN} \label{bab:2}
\petunjuk{Dari kebutuhan mitra ke rancangan solusi. Pada seminar proposal berupa draf yang cukup matang untuk diuji. Analisis kebutuhan dan arsitektur umum boleh sama antaranggota; rancangan bagian Anda diuraikan rinci dan ditulis sendiri. \tahap{\wajib{WAJIB (draf)}}{\wajib{WAJIB (final)}}}

\section{Analisis Kebutuhan} \label{sec:2.1}
\petunjuk{Kebutuhan fungsional dan nonfungsional yang berasal dari mitra, beserta prioritasnya.}
\tulisuraian

\section{Rancangan Solusi} \label{sec:2.2}
\petunjuk{Arsitektur dan rancangan solusi; sisipkan diagram sebagai gambar bernomor dan jelaskan dalam paragraf.}
\tulisuraian

\section{Bukti Kelayakan Awal} \label{sec:2.3}
\petunjuk{Bukti bahwa rancangan layak dikerjakan, misalnya prototipe awal, uji konsep, atau data awal. \tahap{\wajib{WAJIB}}{\wajib{WAJIB (final)}}}
\tulisuraian

\section{Rencana Pengujian dan TKT 3} \label{sec:2.4}
\petunjuk{Untuk setiap uji yang direncanakan, uraikan tujuan uji, kriteria lulus, dan bukti yang akan dikumpulkan menuju TKT minimal 3.}
\tulisuraian
